% !PRINTED 54
\weberpage{54}

\gesperrt{Die Anzahl der Geschlechter ist also so gross, wie die Anzahl der zweiseitigen Elemente}\footnote{\gesperrt{Gauss} hat in die Theorie der quadratischen Formen diese Begriffe zuerst eingeführt, und den Ausdruck \glqq Genera\grqq{} gebraucht. Disqu. arithm. art. 228 f.}.

\begin{center}
\S.\,14.
\end{center}

\begin{center}
\gesperrt{Indices nach einer ungeraden Primzahlpotenz}\\
\gesperrt{als Modul.}
\end{center}

Das wichtigste Beispiel einer commutativen Gruppe bieten die natürlichen Zahlen, wenn sie durch die gewöhnliche Multiplication mit einander verbunden werden.

Um daraus eine endliche Gruppe abzuleiten, nehmen wir eine beliebige ganze positive Zahl $m$ als Modul an und zerlegen $m$ in seine Primfactoren

\weq{1}{m = 2^\lambda q_1^{\varkappa_1} q_2^{\varkappa_2} \ldots,}

worin $q_1, q_2 \ldots$ verschiedene ungerade Primzahlen, $\lambda, \varkappa_1, \varkappa_2 \ldots$ positive Exponenten, $\lambda$ möglicherweise auch die Null, nämlich wenn $m$ ungerade ist, bedeuten.

Nun werden alle Zahlen, positive sowohl als negative, die nach dem Modul $m$ unter einander congruent sind, in eine \gesperrt{Zahlclasse} vereinigt, und jede dieser Zahlclassen wird durch einen Repräsentanten, etwa durch den Rest der Division, also durch eine der Zahlen $0, 1, 2 \ldots m-1$, dargestellt.

Sind $a$ und $a'$ zwei Zahlen einer solchen Classe und $b$ und $b'$ zwei Zahlen einer zweiten Classe, so gehören auch die Producte $ab$ und $a'b'$ in dieselbe Classe. Denn ist $a \equiv a'$, $b \equiv b'$, so ist auch $ab \equiv a'b'\ (\mathrm{mod}\ m)$. Durch die Multiplication werden also nicht bloss die Zahlen, sondern auch die Classen componirt.

Trotzdem bilden diese Zahlclassen in ihrer Gesammtheit noch keine Gruppe; denn aus einer Congruenz

\[
ab \equiv ac\ (\mathrm{mod}\ m)
\]

folgt nur dann nothwendig $b \equiv c$, wenn $a$ relativ prim zu $m$ ist. Es ist also die Forderung \S.\,1, 3. hier im Allgemeinen nicht erfüllt.

Der grösste gemeinschaftliche Theiler, den eine Zahl $a$ mit $m$ hat, ist bei der ganzen durch $a$ repräsentirten Classe der-
%CONT>>
